% !TeX program = lualatex
% Compile twice: lualatex 78.tex
% Standalone research notebook; original section preserved.
% Research additions: 27 September 2026. No external TeX input or BibTeX file.
\documentclass[11pt,a4paper]{article}
\usepackage[margin=24mm,headheight=14pt]{geometry}
\usepackage{fontspec}
\setmainfont{Latin Modern Roman}
\setsansfont{Latin Modern Sans}
\setmonofont{Latin Modern Mono}
\usepackage{amsmath,amssymb,mathtools}
\usepackage{unicode-math}
\setmathfont{Latin Modern Math}
\usepackage{microtype}
\usepackage{enumitem,xurl,needspace}
\usepackage[dvipsnames]{xcolor}
\usepackage{hyperref}
\hypersetup{unicode=true,colorlinks=true,linkcolor=MidnightBlue,urlcolor=MidnightBlue,
 pdftitle={Research notebook - Problem 78},pdfauthor={Open Problems Math},bookmarksnumbered=true}
\usepackage{bookmark}
\usepackage{titlesec}
\titleformat{\section}{\Large\bfseries\raggedright}{\thesection}{0.7em}{}
\usepackage{fancyhdr}
\pagestyle{fancy}\fancyhf{}
\fancyhead[L]{\small\sffamily Open Problems Math | Research notebook}
\fancyhead[R]{\small\sffamily Problem 78 | 27 September 2026}
\fancyfoot[C]{\thepage}
\setlength{\parindent}{0pt}
\setlength{\parskip}{5pt plus 1pt}
\setlength{\emergencystretch}{3em}
\setcounter{tocdepth}{1}
\setcounter{secnumdepth}{1}
\setlist[itemize]{leftmargin=3em,itemsep=5pt,topsep=4pt,parsep=0pt}
\allowdisplaybreaks
\numberwithin{equation}{section}
\newcommand{\N}{\mathbb{N}}
\newcommand{\Z}{\mathbb{Z}}
\newcommand{\Q}{\mathbb{Q}}
\newcommand{\R}{\mathbb{R}}
\newcommand{\C}{\mathbb{C}}
\newcommand{\F}{\mathbb{F}}
\newcommand{\A}{\mathbb{A}}
\newcommand{\PP}{\mathbb P}
\newcommand{\E}{\mathbb{E}}
\newcommand{\eps}{\varepsilon}
\newcommand{\dd}{\,\mathrm d}
\newcommand{\sref}[2]{\hyperlink{source:#1:#2}{[S#2]}}
\newcommand{\eref}[2]{\hyperlink{extra:#1:#2}{[E#2]}}
\newif\ifshowcatalogue
\showcataloguetrue
\newcommand{\cataloguescope}[1]{\ifshowcatalogue
 \paragraph{Exact scope in the supplied catalogue.}
 \begin{quote}\small #1\end{quote}\fi}
\begin{document}
\setcounter{section}{77}
\section{\texorpdfstring{Abundance Conjecture for the minimal model program}{Abundance Conjecture for the minimal model program}}\label{78}
\subsection{Definitions and mathematical statement}
For a projective klt pair in the source's minimal-model setting, put $D=K_X+\Delta$. A divisor is nef if $D\cdot C\ge0$ for every curve $C$. In the usual rational-divisor formulation it is semiample if for some integer $m>0$, $mD$ is Cartier and its line bundle is generated by global sections. The conjecture is
\[
 D\text{ nef}\quad\Longrightarrow\quad \exists m>0:\mathcal O_X(mD)\text{ globally generated}.
\]
The rational-boundary convention is the standard one for this displayed statement; a real-boundary version requires the source's corresponding notion of semiampleness. No bigness hypothesis is added.
\par\smallskip\textit{Formulation and concept references:} \sref{78}{1}.
\cataloguescope{Prove that for every minimal projective Kawamata log terminal pair (X,Delta), the nef adjoint divisor K\_X+Delta is semiample.}
\subsection{Short English statement}
Must a nef canonical-plus-boundary divisor on a minimal klt variety eventually have enough sections to define a morphism?
% BEGIN RESEARCH_ADDITION ATTEMPT_78
\subsection{Research attempt}

\paragraph{Current frontier.}
Lazi\'c studies criteria involving metrics with minimal singularities and
multiplier ideals, including an equivalence under a nonzero Euler
characteristic hypothesis; these additional assumptions are not universal
abundance \eref{78}{1}. Results about anti-canonical nonvanishing in
\sref{78}{4} concern a different divisor. Here the attempt uses the
rational-boundary convention of the displayed target and makes no bigness
assumption. The aim is to isolate the difference between a nef numerical
class, nonzero actual sections, and a section ring whose base ideals
ultimately disappear.

\paragraph{Attack route.}
A metric approach tries to remove singularities, but even a smooth flat
metric can coexist with no sections. A finite-generation approach could
turn asymptotic vanishing of actual base orders into semiampleness, provided
nonvanishing has first been secured. The latter route is pursued by an
explicit base-ideal calculation. The missing geometric steps are then
separated from an invalid replacement of actual base orders by numerical
orders after an ample perturbation.

\paragraph{Attempt: a base-ideal criterion with all hypotheses exposed.}
Let $L$ be a Cartier divisor on a normal projective variety $X$. Assume
there is a positive integer $m_0$ such that
$H^0(X,m_0L)\ne0$ and the Veronese section ring
\[
 R^{(m_0)}(L)=\bigoplus_{k\geq0}H^0(X,km_0L)
\]
is generated by its degree-one piece. Finite generation of the section
ring, together with positive-degree nonvanishing, permits such a sufficiently
divisible Veronese; one can alternatively retain generation in degree one
as the explicit sufficient hypothesis. Let $\mathfrak b$ be the base
ideal of $|m_0L|$, defined by the image of its evaluation map after tensoring
by $\mathcal O_X(-m_0L)$.

Surjectivity of the multiplication maps from the symmetric powers of
$H^0(X,m_0L)$ gives the exact equality of ideal sheaves
\begin{equation}\label{eq78:powers}
 \mathfrak b_{|km_0L|}=\mathfrak b^k\qquad(k\geq1).
\end{equation}
Locally trivialize $m_0L$. Products of $k$ local generators of
$\mathfrak b$ are evaluations of products of global sections, giving one
inclusion; generation of all degree-$k$ sections by sums of such products
gives the other. Thus \eqref{eq78:powers} concerns actual ideals, not just
set-theoretic base loci.

For a divisorial valuation $v$ of the function field with centre on $X$,
write $v(\mathfrak b)$ for the minimum valuation of local generators at
its centre. Equation \eqref{eq78:powers} gives
\begin{equation}\label{eq78:order}
 \frac{1}{km_0}v(\mathfrak b_{|km_0L|})
       =\frac1{m_0}v(\mathfrak b).
\end{equation}
If $\mathfrak b\ne\mathcal O_X$, some such $v$ has
$v(\mathfrak b)>0$. To see this geometrically, blow up the nonzero ideal
$\mathfrak b$ and normalize. Its pullback is an invertible ideal
$\mathcal O(-E)$ with $E$ effective. It is not everywhere a unit: over a
point of the base locus all local generators lie in the maximal ideal at
points above it. Hence $E$ is nonzero. A prime component of $E$ defines
the required valuation. If the original ideal is already invertible, its
nonzero effective Cartier divisor gives the same conclusion directly.

It follows that the simultaneous hypotheses
\begin{equation}\label{eq78:criterion}
 \begin{gathered}
 H^0(X,m_0L)\ne0,\qquad R^{(m_0)}(L)\text{ generated in degree one},\\
 \lim_{k\to\infty}\frac{v(\mathfrak b_{|km_0L|})}{km_0}=0
       \quad\text{for every divisorial }v.
 \end{gathered}
\end{equation}
force $\mathfrak b=\mathcal O_X$. Thus $m_0L$ is globally generated and
$L$ is semiample. The proof uses neither a numerical approximation to
$v$ nor an exchange of a limit with a minimum over all valuations: one
fixed positive valuation would contradict the last condition.

\paragraph{Attempt: application to the adjoint target.}
For the original rational pair choose $q>0$ with
$L=q(K_X+\Delta)$ Cartier. If the nonvanishing, generation and actual
base-order assertions in \eqref{eq78:criterion} can be established for
this $L$, the target follows, since a globally generated multiple of $L$
is a globally generated multiple of $K_X+\Delta$.
This packages three distinct arithmetic-geometric obligations. It does
not assert that nefness alone proves any of them. In particular, a section
ring with no positive-degree sections is trivially finitely generated but
useless for semiampleness; the first condition excludes this pitfall.

An analytic continuation of the route might compare the actual base ideals
with multiplier ideals of minimal metrics, then use vanishing or Euler
characteristic information to produce sections. That is exactly where
extra hypotheses in the metric literature matter \eref{78}{1}.
No equality between those metric ideals and the actual base ideals is
proved here. Without it, substituting metric Lelong numbers for the
valuations in \eqref{eq78:criterion} is an unproved bridge.

\paragraph{Stress test: a nef flat line bundle with no sections.}
Let $C$ be a complex elliptic curve and choose a non-torsion degree-zero
line bundle $P\in\operatorname{Pic}^0(C)$. Such bundles exist because the
complex torus $\operatorname{Pic}^0(C)$ is uncountable and its torsion
subgroup is countable. For every $m>0$, a nonzero section of $P^m$ would
have an effective divisor of degree zero, hence no zeros, and would
trivialize $P^m$. This contradicts non-torsion. Therefore
\[
 P\text{ is nef},\qquad H^0(C,P^m)=0\quad(m>0).
\]
It admits a smooth unitary flat metric; its curvature and Lelong numbers
are zero. Also $P+\varepsilon A$ has positive degree for an ample $A$
and $\varepsilon>0$, so numerical vanishing orders after ample perturbation
carry no obstruction of this sort. These facts show why numerical
positivity and zero metric singularities alone do not create sections.

This is not an abundance counterexample. On an elliptic curve an effective
boundary $\Delta$ with $\deg(K_C+\Delta)=0$ must be zero, and $K_C$ is
trivial. The non-torsion $P$ cannot be the stipulated adjoint divisor in
that degree-zero situation. The missing adjoint input is therefore genuine,
not a reason to discard the conjecture.

There is another weak argument to avoid. Along factorial multiples, base
loci form a descending family of closed sets whenever the relevant sections
are present. Noetherianity can make that family stabilize at a nonempty
closed set; it does not make the set disappear. Likewise finding more
sections at each step does not imply a section avoiding every remaining
base point. The valuation in \eqref{eq78:order} detects such a persistent
base ideal quantitatively only after the generation hypothesis has been
established.

\paragraph{Outcome.}
\textbf{Outcome: Conditional reduction.}

A direct base-ideal proof reduces semiampleness to nonvanishing,
degree-one generation after a Veronese, and zero actual asymptotic base
orders. It isolates an invalid metric-to-section shortcut by an explicit
nef flat example outside the adjoint class. The argument supplies no new
unconditional abundance case and makes no priority claim for its elementary
base-ideal lemma.

The unresolved task is to prove the three properties for arbitrary nef
klt adjoint divisors without adding bigness, Euler characteristic, or a
lower-dimensional MMP assumption. A real-boundary formulation would also
need its own notion of semiampleness; it is not proved by this
rational-multiple calculation.
% END RESEARCH_ADDITION ATTEMPT_78
\subsection{Sources}
\begin{itemize}\raggedright
\item[\textnormal{[S1]}]\hypertarget{source:78:1}{} Vladimir Lazic, Metrics with minimal singularities and the Abundance conjecture (2024).
\url{https://arxiv.org/abs/2406.18233}.
{\small\textit{Catalogue formulation source.}}
\item[\textnormal{[S2]}]\hypertarget{source:78:2}{} Inequalities of Miyaoka-Yau type and Uniformisation of varieties of intermediate Kodaira Dimension.
\url{https://arxiv.org/abs/2601.15138}.
{\small\textit{Catalogue research/status reference; not a proof endorsement.}}
\item[\textnormal{[S3]}]\hypertarget{source:78:3}{} An approach to the abundance conjecture for Kaehler varieties via algebraic reduction.
\url{https://arxiv.org/abs/2604.03879}.
{\small\textit{Catalogue research/status reference; not a proof endorsement.}}
\item[\textnormal{[S4]}]\hypertarget{source:78:4}{} Two non-vanishing results concerning the anti-canonical bundle.
\url{https://www.cambridge.org/core/journals/nagoya-mathematical-journal/article/two-nonvanishing-results-concerning-the-anticanonical-bundle/1E68F2B88B818A7CBD47B96B7D700548}.
{\small\textit{Catalogue research/status reference; not a proof endorsement.}}
% BEGIN RESEARCH_ADDITION REFERENCES_78
\item[\textnormal{[E1]}]\hypertarget{extra:78:1}{} Vladimir Lazi\'c,
\emph{Metrics with minimal singularities and the Abundance conjecture},
arXiv:2406.18233, introduction and hypotheses of the metric/nonvanishing
criteria, particularly the nonzero Euler characteristic setting.
\url{https://arxiv.org/abs/2406.18233}.
{\small\textit{Inspected 27 September 2026; the notebook does not assume
that metric singularity control automatically identifies all actual base ideals.}}
% END RESEARCH_ADDITION REFERENCES_78
\end{itemize}

\end{document}
